\documentclass[12pt]{article}
\usepackage[utf8]{inputenc}
\usepackage[T1]{fontenc}
\usepackage{hyperref}
\title{Slavomolisano Language Resources for K-12 Curriculum Development}
\author{Compiled by Claude Sonnet 5.5 (claude-sonnet-5-5) for Andr\'as Kornai}
\date{}
\begin{document}
\maketitle

\section{Introduction}
Slavomolisano (Molise Croatian, ISO 639-3 svm) is a Shtokavian Croatian variety spoken in three villages of Campobasso province, Italy: Montemitro, Acquaviva Collecroce and San Felice del Molise. One source reports fewer than 1,000 active speakers and fewer than 2,000 passive speakers. Italian dominates administration, church and school. I found no school programme for it in my searches.

\section{Primary Educational Resources}
I found no primers, textbooks or school programmes.

\section{Linguistic and Reference Materials}
\begin{itemize}
\item Walter Breu, \emph{Grammatica essenziale dello slavomolisano}, an open-access (CC BY-SA 4.0) basic grammar hosted at the University of Konstanz repository: \url{https://kops.uni-konstanz.de/server/api/core/bitstreams/32ad082a-52a2-4da0-8869-dd618d3a17d9/content}
\item The Konstanz Slavic linguistics page on the language, associated with the Acquaviva Collecroce dictionary (\emph{Dizionario croato molisano di Acquaviva Collecroce}, Breu and Piccoli, 2000): \url{https://uni-konstanz.de/FuF/Philo/Sprachwiss/slavistik/acqua/spracheE.htm}
\item Slavomolisano (Wikipedia): \url{https://wikipedia.com/wiki/ISO_639:svm}
\item Verbix documentation on Slavomolisano: \url{https://docs.verbix.com/Languages/Slavomolisano}
\end{itemize}

\section{Religious and Community Resources}
I found nothing in my searches.

\section{Digital and Audio Resources}
\begin{itemize}
\item A University of Salerno repository record that came up in the search. I did not examine its contents: \url{https://www.iris.unisa.it/handle/11386/4819215}
\end{itemize}

\section{Curriculum Development Context}
The language is used mainly by older adults and is not known to be taught in schools. The open grammar and the Breu--Piccoli dictionary are the best starting points for any teaching material. A curriculum would have to begin with informal or extracurricular learning and with the standardisation based on Acquaviva usage.

\end{document}